\scanpage{11}
$\Omega$ module de différentielles pour un C.M. $\Longleftrightarrow$ les $\Omega/(f_1\ldots f_n)\Omega$ sont les injectifs indécomposables

\deleted{\uline{Lemme \unsure{1}} Soit $X$ un préschéma régulier, \unsure{un} \illegible{1} \unsure{fibré}} \deleted{\unsure{fini}}

\uline{Lemme 8} Soit $X$ un préschéma de Cohen-Macaulay ; \margin{\tikz{\draw (0,0)--(0,1.8); \draw (.08,0)--(.08,1.8);}} $\Omega$ un faisceau fondamental sur $X$, alors pour tte suite $(f_1\ldots f_n)$ de sections de $\mathcal O_X$ qui est $\mathcal O_X$-régulière, $\Omega/(f_1,\ldots f_n)\Omega$ est un faisceau fondamental sur le \deleted{\unsure{pré}}schéma \deleted{\unsure{noethérien|noeth.}} \inserted{CM} \deleted{$\unsure{H^0(|H^2(}$} $V(f_1,\ldots f_n)$.

\uline{Corollaire} 9. Soit $x\in X$ et $f_1\ldots f_n\in\mathcal O_{X,x}$ un syst. de paramètres de $\mathcal O_{X,x}$, alors \margin{\tikz{\draw (0,0)--(0,1.7);}} $\Omega/(f_1\ldots f_n)\Omega$ est \deleted{\unsure{un}} l'injectif indécomposable dans la catégorie des modules sur $\mathcal O_{X,x}/(f_1\ldots f_n)$.

\uline{Corollaire 10} Soit $X$ un préschéma loc. noeth \margin{\tikz{\draw (0,0)--(0,1.5);}} qui est \uline{localemt immergeable dans} un \uline{régulier}.
% slanted
\margin{il suffit \deleted{\unsure{qu'il|que}} qu'il soit excellent, ou seulement qu'il satisfasse les conditions de EGA \uline{IV} 6.11.8.}
Alors il existe un ouvert dense $U$ \deleted{de} $X$, et un Module coh. $\Omega$ sur $X$, ayant les propriétés suivantes

a) Les anneaux locaux de $U$ sont C.M.

b) Pour tout $x\in U$, et tt syst. de paramètres $f_1\ldots f_n$ de $\mathcal O_{X,x}$, le module $\Omega/(f_1\ldots f_n)\Omega$ sur $\mathcal O_{X,x}/(f_1,\ldots f_n)$ est

% archive sheet number: 10
